\section*{الفصل \ref{ch:g2:where-animals-live} --- أين تعيش الحيوانات}

\begin{solution}{exo:g2:where-animals-live:1}
نوع المكان الذي يعيش فيه \omterm{def:g1:animals-around-us:animal}{الحيوان} على نحو طبيعي. وهو يجد فيه غذاءه
وماءه ومأواه وموضعًا آمنًا لصغاره.
\end{solution}

\begin{solution}{exo:g2:where-animals-live:2}
البركة: ضفدع، ورعّاشة (أو سمكة شوكية، أو حلزون ماء). الغابة:
سنجاب، ونقّار خشب (أو أيّل، أو نملة خشب). المرج: أرنب، وجندب (أو
فراشة، أو فأر حقل).
\end{solution}

\begin{solution}{exo:g2:where-animals-live:3}
السمكة الشوكية: البركة. نقّار الخشب: الغابة. الجندب: المرج. حشرة
القمل الخشبي: تحت حجارة جدار قديم (الزوايا الرطبة المختفية).
\end{solution}

\begin{solution}{exo:g2:where-animals-live:4}
الضفدع: قدمان مكفوفتان تدفعان الماء. السنجاب: مخالب قابضة وذنب
كبير موازن.
\end{solution}

\begin{solution}{exo:g2:where-animals-live:5}
يدان مجرفتان عريضتان للحفر؛ \omterm{def:g1:animals-around-us:covering}{وفرو} ينبطح في الاتجاهين فلا تشعّثه
التربة؛ \omterm{def:g1:the-five-senses:sense}{وعينان} لا يكاد يحتاج إليهما في الظلام.
\end{solution}

\begin{solution}{exo:g2:where-animals-live:6}
نقّار الخشب: عاليًا على الجذوع. السنجاب: في \omterm{def:g1:plants-around-us:tree}{الأغصان}. الأيّل: بين
الشجيرات. نملة الخشب: على الأرض، في عشها المكوّم.
\end{solution}

\begin{solution}{exo:g2:where-animals-live:7}
تختفي مع البركة --- فمن غير \omterm{def:g2:where-animals-live:habitat}{الموئل} تفقد الغذاء والمأوى ومكان
التوالد دفعةً واحدة. ولذلك تبدأ حماية \omterm{def:g1:animals-around-us:animal}{الحيوانات} بحماية الأماكن
التي تعيش فيها.
\end{solution}

\begin{solution}{exo:g2:where-animals-live:8}
الجواب مفتوح. وحتى حجر واحد يُقلب كثيرًا ما يؤوي حشرات قمل خشبي،
ونملًا، وخنفساء، وأحيانًا أم أربعة وأربعين --- \omterm{def:g2:where-animals-live:habitat}{فالموائل} الصغيرة
ملأى على نحو مدهش.
\end{solution}

\begin{solution}{exo:g2:where-animals-live:9}
قدمان مكفوفتان (للتجديف)، \omterm{def:g1:animals-around-us:covering}{وريش} يدفع الماء، ومنقار للغربلة: وكلها
تدل على الماء --- \omterm{def:g2:where-animals-live:habitat}{موئل} بركة أو بحيرة. \omterm{def:g1:animals-around-us:animal}{والحيوان} بطة على الأرجح.
\end{solution}

\begin{solution}{exo:g2:where-animals-live:10}
عتاد الضفدع --- القدمان المكفوفتان، \omterm{def:g1:the-five-senses:sense}{والجلد} الرطب، وجسم مصنوع
للماء وللضفاف الندية --- لا فائدة منه في صندوق جاف، وموئله أعطاه
أكثر من الغذاء: ماءً لجلده ولبيضه، ومخابئ، وضفادع أخرى. والإحسان
إلى \omterm{def:g1:animals-around-us:animal}{حيوان} بري هو أن يُترك في موئله --- أو يُعاد إليه.
\end{solution}
